%----------------------------------------------------------------------
\subsection{Session 12: the minorisation route is closed by inert dust;
(C2) is polynomial, not geometric}\label{sec:s12}

Session 12 was directed at (C1)+(C2) of \S\ref{sec:s11left}, through
TODO~12a (a Lyapunov drift for $\lambda$), 12b (the minorisation of
Proposition~\ref{prop:c1}) and 12c (an odd perturbation), with the
instruction to return one of (A1)--(A3) or a clean negative (B).  The
outcome is (B), in a form sharper than the one the brief anticipated.
Four findings, in logical order.

\begin{enumerate}[leftmargin=2.2em]
\item[(1)] \textbf{The environment carries inert dust, and this closes
every Doeblin-type route.}  A curve of mass $\varepsilon$ is touched by
a chord with probability $\le2\varepsilon$ per step, so it survives
\emph{unchanged} through the first $n$ flips with probability
$\to1$ as $\varepsilon\to0$ (Lemma~\ref{lem:inert}).  Consequently the
post-flip laws from the states $x_\varepsilon$ (a fixed macroscopic
marked pair, plus one dust curve of mass $\varepsilon$) are supported,
up to $o(1)$, on the pairwise disjoint sets ``some curve has mass
exactly $\varepsilon$''.  No common minorant exists: the hypothesis of
Proposition~\ref{prop:c1} fails on $\Omega_K$ for \emph{every} $n$
(Proposition~\ref{prop:nomino}), and so does uniform ergodicity of
$\kappa_1$ in $L^\infty(\Omega_K)$ and $V$-uniform ergodicity for
\emph{every} weight $V$ (Corollary~\ref{cor:noerg}).  Route (a) of
\S\ref{sec:s11left} and TODO~12b are therefore closed, not
deferred.  The proof uses only $\tau_n<\infty$ a.s.
\item[(2)] \textbf{TODO~12a fails on the full state space.}  From a
state whose non-dust material has total mass $m$ and whose remaining
mass sits in $N\gg m^{-3}$ equal pieces, one has
$\sum_jr_j^4\le17m^4$ deterministically for a long time, so
$\kappa_1^kV\ge\tfrac34(17m^4)^{-\theta}$ while $V\le m^{-4\theta}$:
no inequality $\kappa_1V\le\varrho V+b$ with $\varrho<1$ can hold
(Proposition~\ref{prop:nodrift}).  As instructed, the fallback is the
killed chain on $\Omega_K$, where $V$ is bounded anyway.
\item[(3)] \textbf{(C2) is true but polynomial.}  With $2\times10^7$
samples from the one-curve start,
$a_k(x)=\E_x[(-1)^{\tau_k}]=(-1)^k\,(2.5\pm0.1)\,k^{-3}$ for
$7\le k\le13$, and with $10^8$ samples the local exponent over
$8\le k\le20$ is $2.7\pm0.3$; a tail of the same form,
$k^3|a_k|\to2.0$--$2.4$, from the three-curve start $(0.3,0.3;\{0.4\})$,
and the same tail \emph{emerging} after a transient of $\approx8$
flips from a nine-curve start.  Session~11's ``$\rho\approx0.3$''
was the transient above a noise floor of $0.003$.  Since
$a_k=\pm\E_x[(-1)^{B(\bar\Gamma_{\tau_k})}]$ is the parity bias of the
block count at the $k$-th flip, the mechanism is the slow ($\log$)
growth of $B$: the block-count parity is a slow mode of $\kappa_1$ at
eigenvalue $+1$, not a mode near $-1$ (\S\ref{sec:s12num}).  The
polynomial decay is insensitive to inserting dust into the start
(identical to two standard errors) and is \emph{weaker} from starts
with $9$ or $17$ curves.
\item[(4)] \textbf{What the reduction actually needs is polynomial.}
At $z=-1$ Proposition~\ref{prop:condecay} needs only
$G\in C^{0,\gamma}$, $\gamma>4/9$, and a polynomial decay of the
twisted parity $\E_x[w^{\tau_k}(-1)^{E_{\tau_k}}]$ with exponent $p$,
uniformly for $w$ near $1$, gives $\gamma=(p-1)/(p+1)$ by a crude
bound (Proposition~\ref{prop:holder}), so the budget needs $p>13/5$;
if the parity memory factorises from the time weighting, as it does
numerically, the exponent improves to $\min(1,p-1)$ and the budget
needs only $p>13/9$.  The measured $p=2.7\pm0.3$ meets the second
comfortably and the first marginally.  (C2) is accordingly replaced by
(C2$''$) in \S\ref{sec:s12left}.  Any proof must be \emph{dust-blind},
i.e.\ work in a topology in which $x$ and $x_\varepsilon$ are close;
we record why the natural transport metric is not contracted by the
dynamics and why an adversarial environment cannot be allowed
(\S\ref{sec:s12left}).
\end{enumerate}

\subsubsection{Inert dust}\label{sec:s12dust}

Throughout, the chain is realised on the circle: the state is a
partition of the circle into curves (each a cyclically ordered union of
arcs, the cyclic order being that of the chord diagram), a step picks
two independent uniform points, and the move is determined by the
curves the points lie in and their positions along those curves.

\begin{lemma}[inert dust coupling]\label{lem:inert}
Let $y$ be a state, $D$ a sub-arc of the circle of mass $\varepsilon$
lying inside one curve of $y$ and containing neither mark, and let
$y_\varepsilon$ be the state obtained from $y$ by declaring $D$ a
separate (unmarked) curve.  Drive the chains from $y$ and from
$y_\varepsilon$ by the same uniform points.  Then, up to the first step
at which a point falls in $D$, the state of the $y_\varepsilon$-chain
is the state of the $y$-chain with $D$ declared a separate curve; in
particular the two chains have the same flip times, and the
$y_\varepsilon$-chain contains a curve of mass exactly $\varepsilon$.
The first point in $D$ falls at a step of probability at most
$2\varepsilon$, so for every $T$,
\[
\Pr\bigl[\text{the coupling holds through time }T\bigr]\ge1-2\varepsilon T .
\]
\end{lemma}

\begin{proof}
$D$ is a contiguous segment of the cyclic order of the curve $C\supset D$
containing it.  A split of $C$ at two points outside $D$ cuts the cycle
into two arcs, one of which contains $D$; in the $y_\varepsilon$-chain
the same cuts split $C\setminus D$ into the same two arcs with $D$
excised.  A merge of $C$ with another curve at points outside $D$
splices the two cycles at those points; in the $y_\varepsilon$-chain the
same splice is applied to $C\setminus D$.  Moves not involving $C$ are
identical.  A flip is a split or merge decided by which marked material
the two points lie in; the marked material of the two chains differs
only by $D$, which contains no point.  Induction on the steps.
\end{proof}

\begin{proposition}[no state-uniform minorisation]\label{prop:nomino}
Let $K\ge4$ and assume $\tau_n<\infty$ a.s.\ from the state
$x_0=(\tfrac14,\tfrac14;\{\tfrac12\})$ (apart sector).  Then there are
no $n\ge1$, $\delta>0$ and probability measure $\nu$ on $\Omega_K$ with
\[
\Pr^{\Omega_K}_{\bar x}\bigl[\bar\Gamma_{\tau_n}\in\cdot\,\bigr]\ge\delta\,\nu(\cdot)
\qquad\text{for every }\bar x\in\Omega_K\cap S'_K ,
\]
for the chain killed on $\Omega_K$ or for the unkilled chain.  A
fortiori the hypothesis of Proposition~\ref{prop:c1} is never
satisfied.
\end{proposition}

\begin{proof}
Suppose such $n,\delta,\nu$ exist.  For $i\ge3$ let
$\varepsilon_i=2^{-i}$, let $x_i$ be $x_0$ with a sub-arc $D_i$ of mass
$\varepsilon_i$ of its environment curve declared a separate curve
(so $x_i=(\tfrac14,\tfrac14;\{\tfrac12-\varepsilon_i,\varepsilon_i\})$,
$B(x_i)=4$, $x_i\in\Omega_K\cap S'_K$), and let
$S_i=\{\bar x\in\Omega_K:\ \text{some curve of }\bar x\text{ has mass
exactly }\varepsilon_i\}$.  Apply the minorisation at $x_i$ to the set
$\Omega_K\setminus S_i$.  By Lemma~\ref{lem:inert} with $T=T_i$, on the
event that the coupling holds through $T_i$ and $\tau_n\le T_i$ for the
unkilled $x_0$-chain, the $x_i$-chain performs its $n$-th flip at the
same time and, if it is still alive, its state then contains $D_i$
intact; if it has been killed, its state is the cemetery, which is not
in $\Omega_K\setminus S_i$.  Hence
\[
\delta\,\nu(\Omega_K\setminus S_i)\;\le\;2\varepsilon_iT_i+\Pr_{x_0}[\tau_n>T_i]
\qquad\text{for every }T_i ,
\]
and $T_i=\varepsilon_i^{-1/2}$ gives $\nu(S_i)\to1$ as $i\to\infty$.
But a configuration in $\Omega_K$ has at most $K$ curves, hence at most
$K$ distinct masses, hence lies in at most $K$ of the sets $S_i$; so
$\sum_i\nu(S_i)=\int\#\{i:\bar x\in S_i\}\,\nu(d\bar x)\le K<\infty$,
contradicting $\nu(S_i)\to1$.
\end{proof}

\emph{Status: PROVED.  Note what the proof does not use: no moment of
$\tau$, no property of the environment beyond the trivial one that a
curve of mass $\varepsilon$ is rarely hit, and no bound on $K$ beyond
$K\ge4$.  It also shows that the difficulty is not the one anticipated
in TODO~12b (``$K^{-2K}\to K^{-O(1)}$''): the obstruction is not the
\emph{cost} of driving the environment to a reference configuration but
the fact that no finite sequence of steps can remove a curve that is
never hit.  The same applies with the $(n,n+1)$ pair of
Proposition~\ref{prop:c1}, to any weighted space, and to any
``small set'' containing a state with a dust curve --- which every set
of positive $\bar\pi$-measure does, since dust is created at every
scale at rate $\asymp2\eta\,d\eta$ per step.}

\begin{corollary}[no uniform or $V$-uniform ergodicity]\label{cor:noerg}
Let $K\ge4$ and let $\kappa_1$ be the post-flip kernel of the unkilled
chain.
\begin{enumerate}[leftmargin=2em]
\item[(i)] There are no probability measure $\bar\pi$, $C<\infty$ and
$\rho<1$ with $\sup_{\bar x\in\Omega_K}|\kappa_1^kF(\bar x)-\bar\pi F|
\le C\rho^k\|F\|_\infty$ for all bounded $F$ and all $k$.
\item[(ii)] Let $\mathcal Y=\{(a',\tfrac12-a';\{\tfrac12\}):
a'\in[\tfrac18,\tfrac38]\}$ and assume
$(\mathrm M_p)$: $\sup_{y\in\mathcal Y}\E_y[\tau_k]\le Ck^p$ for some
$C,p<\infty$.  Then there is no $V\ge1$ for which $\kappa_1$ is
$V$-uniformly ergodic: no $\bar\pi$, $C'$, $\rho<1$ with
$|\kappa_1^kF(\bar x)-\bar\pi F|\le C'\rho^kV(\bar x)\|F\|_V$ for all
$F$ with $\|F\|_V=\sup|F|/V<\infty$ and all $\bar x$ with $B(\bar x)\le4$.
\end{enumerate}
\end{corollary}

\begin{proof}
(i) With $F_i=\mathbf 1_{S_i}$ (notation of the previous proof) and
$x_i$ as there, Lemma~\ref{lem:inert} gives
$\kappa_1^kF_i(x_i)\ge1-2\varepsilon_iT-\Pr_{x_0}[\tau_k>T]$, while
$\kappa_1^kF_i(x_0)=0$ because from $x_0$ every curve mass at every time
is either a sum of some of $\tfrac14,\tfrac14,\tfrac12$ or has an
atomless law, so no curve ever has mass exactly $\varepsilon_i\le\tfrac18$ a.s.  Let
$i\to\infty$: the oscillation of $\kappa_1^kF_i$ over $\Omega_K$ tends
to $1$ for each fixed $k$, contradicting $2C\rho^k<1$.

(ii) Suppose such $V,\bar\pi,C',\rho$ exist.  First, $\Pi=\mathbf 1\otimes\bar\pi$
is bounded on the $V$-space, so $\bar\pi V<\infty$ and
$\kappa_1V\le(C'\rho+\bar\pi V)\,V$; in particular $\kappa_1V(x_0)<\infty$.
Second, for $y\in\mathcal Y$ and $y_\varepsilon$ obtained by carving a
sub-arc of mass $\varepsilon$ from its environment curve, the argument
of (i) at $y$ gives $\kappa_1^kF_i(y_\varepsilon)\ge1-2\varepsilon T-\Pr_y[\tau_k>T]$
and $\bar\pi F_i=\lim_k\kappa_1^kF_i(x_0)=0$, so with $T=T_k:=4Ck^p$
and $\varepsilon\le\varepsilon_k:=1/(8T_k)$ (Markov and $(\mathrm M_p)$),
\[
V(y_\varepsilon)\;\ge\;\frac{1-\tfrac14-\tfrac14}{C'\rho^k}\;=\;\frac1{2C'\rho^k}
\qquad(y\in\mathcal Y,\ \varepsilon\le\varepsilon_k).
\]
Third, from $x_0$ the first step carves a piece of mass in
$d\varepsilon$ from the environment curve with probability
$\ge\tfrac12\,d\varepsilon$ (pick that curve twice, $\tfrac14$;
$|t_1-t_2|/\tfrac12$ has density $\ge2$ near $0$) and the second step
is a flip landing in $\mathcal Y$ with probability $\ge\tfrac18\cdot\tfrac34$;
the landing state is then some $y_\varepsilon$, $y\in\mathcal Y$.  So
\[
\kappa_1V(x_0)\;\ge\;\frac{3}{64}\sum_{k\ge1}\frac{\varepsilon_k-\varepsilon_{k+1}}{2C'\rho^k}
\;\ge\;c\sum_k\frac{k^{-p-1}}{\rho^{k}}\;=\;\infty ,
\]
contradicting the first step.
\end{proof}

\emph{Status: PROVED ((ii) under $(\mathrm M_p)$, which (R1) along the
chain implies with $p=1$ and which is measured: $\E[\tau_k]\approx6k$).
For the chain killed on $\Omega_K$ the same argument runs up to
$k\le c\,3^{K/(3p)}$, so any $V$-uniform ergodicity of the killed chain
has $(1-\rho)^{-1}\log(C'V(x_0)\bar\pi V)\ge c\,3^{K/(3p)}$: exponential
in $K$, i.e.\ fatal in the budget of the session-11 handoff.  This
settles, negatively, the sentence of \S\ref{sec:markovrenewal} that
asked for a weighted space ``in which the post-flip chain is
$V$-geometrically ergodic in the base coordinates'': there is none.
It does \emph{not} touch (C1) or (C2) themselves, whose null vectors
would have to sit near $-1$, whereas the dust modes sit near $+1$.}

\subsubsection{The drift for \texorpdfstring{$\lambda$}{lambda} fails
on the full state space (TODO~12a)}\label{sec:s12drift}

\begin{proposition}[no Lyapunov drift for $(\sum r_j^4)^{-\theta}$]
\label{prop:nodrift}
Let $V=(\sum_jr_j^4)^{-\theta}$ with $0<\theta\le\tfrac12$, on the full
(unkilled) state space, and assume $\tau_k<\infty$ a.s.\ from every
state.  There are no $\varrho<1$ and $b<\infty$ with
$\kappa_1V\le\varrho V+b$.
\end{proposition}

\begin{proof}
Suppose the inequality holds.  Iterating,
$\kappa_1^kV\le\varrho^kV+b/(1-\varrho)$ for all $k$.  Fix $k$ with
$\varrho^k\le\tfrac38\,17^{-\theta}$ and $m\in(0,\tfrac14)$.

\emph{The isolated process.}  Let $\bar x$ be the apart-sector state
with $a=m$, $b=m/10$ and one environment curve of mass $1-1.1m$, and
let $\mathrm P^{\rm iso}$ be the law of the chain from $\bar x$ modified
so that a chord with one point in the original environment curve (or in
any curve descended from it by splits and merges among such curves)
and one point elsewhere does nothing.  Under $\mathrm P^{\rm iso}$ the
``non-dust'' material of total mass $1.1m$ evolves as a split--merge
chain in its own right, slowed by the factor $(1.1m)^2$, with the
marked pair inside it; its flip times are a.s.\ finite, so there is
$T=T(m,k)$ with $\mathrm P^{\rm iso}[\tau_k>T]\le\tfrac18$.

\emph{The states $x_{m,N}$.}  Let $x_{m,N}$ be $\bar x$ with the
environment curve cut into $N$ equal pieces,
$N\ge16(T+1)^3\max(m^{-1},m^{-4/3})$.  Drive the chain from $x_{m,N}$
and the isolated process from $\bar x$ by the same points (the dust
pieces of $x_{m,N}$ are sub-arcs of the environment curve of $\bar x$;
dust--dust chords split or merge dust material in both and never
touch the non-dust material).  The two processes make the same non-dust moves
except at chords joining dust to non-dust, of which there are at most
$T$ up to time $T$; each such chord appends to one non-dust curve of the
$x_{m,N}$-chain a dust curve, i.e.\ a union of at most $T+1$ original pieces or
fragments of them, of mass $\le(T+1)/N$.  As in Lemma~\ref{lem:inert} (with the
appended material in the role of $D$, now inside a curve of one chain
and absent from the other), the non-dust configurations of the two
processes then agree up to the appended material, and their flip times
agree, until a point falls in appended material, which has probability
$\le2T(T+1)/N$ per step.  Hence
$\Pr_{x_{m,N}}[\tau_k>T]\le\tfrac18+2T^2(T+1)/N\le\tfrac14$.

\emph{The deterministic bound.}  Along the $x_{m,N}$-chain up to time
$T$, the non-dust mass is at most $1.1m+T(T+1)/N\le2m$ and every dust
curve has mass $\le(T+1)/N\le m^{4/3}$, so
\[
\sum_jr_j(\bar\Gamma_t)^4\;\le\;(2m)^4+\Bigl(\max_{\text{dust}}r_j\Bigr)^3
\sum_{\text{dust}}r_j\;\le\;16m^4+m^4=17m^4\qquad(t\le T),
\]
whence $\kappa_1^kV(x_{m,N})\ge\Pr_{x_{m,N}}[\tau_k\le T]\,(17m^4)^{-\theta}
\ge\tfrac34\,17^{-\theta}m^{-4\theta}$, while $V(x_{m,N})\le m^{-4\theta}$
because the $u$-curve alone contributes $m^4$.  The drift bound gives
$\tfrac34\,17^{-\theta}m^{-4\theta}\le\tfrac38\,17^{-\theta}m^{-4\theta}+b/(1-\varrho)$,
i.e.\ $m^{-4\theta}\le\tfrac83\,17^{\theta}b/(1-\varrho)$ for every
$m\in(0,\tfrac14)$: false.
\end{proof}

\emph{Status: PROVED.  This is the fallback the brief provided for: the
obstruction is the family ``all non-dust mass $\asymp m$, the rest in
ultra-fine mass-carrying dust'', which is of vanishing probability
under the chain from one curve (Corollary~\ref{cor:horizon}) but is a
legitimate state, and on it $\lambda$ cannot recover within one
inter-flip period because there is no macroscopic partner to merge
with.  On $\Omega_K$, $V\le(2K)^{4\theta}$ is bounded and the drift is
vacuous.  The killing at $B=K$ in Theorem~\ref{thm:idleloop} therefore
stays, and it is harmless: it costs $(9d+1)^{-\gamma}$ over the horizon.
The second purpose of TODO~12a --- supplying the base half of the
minorisation --- is void by Proposition~\ref{prop:nomino}.  A quick
numerical check of the mechanism (\texttt{s12\_driftfail}): from
$(0.3,0.03;\ 3\times10^4\text{ equal pieces})$, $\E[(\sum r^4(\bar\Gamma_\tau)/\sum r^4(x))^{-\theta}]=1.19$
$(\theta=0.1)$, $1.71$ $(\theta=0.3)$, against $0.91$, $0.77$ with a
single environment curve of the same mass.}

\subsubsection{TODO~12c: the stated criterion does not give (C2)}
\label{sec:s12odd}

The session-11 handoff proposed, as route 12c, a state-uniform lower
bound $c(K)$ on the probability that two consecutive inter-flip times
have opposite parity, ``which gives (C2) by a standard two-block
argument''.  It does not: if the inter-flip times were the
deterministic sequence $1,2,1,2,\dots$ then consecutive parities are
always opposite ($c=1$) while $(-1)^{\tau_k}$ is periodic with period
$4$ and does not decay.  What the two-block argument needs is the
\emph{operator} statement $\|\kappa_{-1}^{\,2}\|_\infty<1$, and that is
false exactly: $\kappa_{-1}(-1)^E=-(-1)^E$ (Proposition~\ref{prop:r2false}),
so $\|\kappa_{-1}^{\,m}\|_\infty=1$ for every $m$.  Any odd perturbation
that is exact returns a state with a different block-count parity,
which is visible; the only odd perturbations that are invisible are
dust insertions, whose weight at scale $\eta$ is $\asymp2\eta$ per step
and whose lifetime is $\asymp1/(2\eta)$, and one checks in one line
that they give a contraction $1-O(\eta)$ per flip and nothing
geometric.  Route 12c is closed as stated; the honest form of (C2) is
in \S\ref{sec:s12num} and \S\ref{sec:s12left}.

\subsubsection{The decay of the parity along the post-flip chain is
polynomial}\label{sec:s12num}

By Lemma~\ref{lem:parity}, $(-1)^{\tau_k}=(-1)^{B(\bar\Gamma_{\tau_k})-B(\bar x)}$
exactly, so
\begin{equation}\label{eq:akB}
a_k(\bar x):=\E_{\bar x}[(-1)^{\tau_k}]=(-1)^{B(\bar x)}\sum_{b\ge1}(-1)^b\,
\Pr_{\bar x}\bigl[B(\bar\Gamma_{\tau_k})=b\bigr],
\end{equation}
the parity bias of the block count at the $k$-th flip, and
$\kappa_1^{\,k}(-1)^E=(-1)^{E}(-1)^ka_k$, so (C2) is the statement that
\eqref{eq:akB} decays geometrically, uniformly in $\bar x$.

\emph{Measurement (\texttt{s12\_akdecay}, \texttt{s12\_akB}; fixed
deterministic starts, $2\times10^7$ independent trajectories each,
standard error $2.2\times10^{-4}$, no censoring).}  From the one-curve
start (arcs $(\tfrac12,\tfrac12)$):
\[
a_k=-0.4003,\ +0.1530,\ -0.0674,\ +0.0335,\ -0.0186,\ +0.0111,\ -0.0073,\
+0.0048,\ -0.0034,\ +0.0022,\ -0.0019,\ +0.0015\quad(k=1,\dots,12),
\]
so that $k^3|a_k|=2.51,2.48,2.49,2.20,2.54,2.54$ for $k=7,\dots,12$:
\[
a_k(\text{one curve})\;=\;(-1)^k\,(2.5\pm0.1)\,k^{-3}\qquad(7\le k\le12),
\]
with the successive ratios $|a_{k+1}/a_k|$ rising through
$0.38,0.44,0.50,0.56,0.60,0.66,0.66,0.70$ --- no geometric rate.  A
$10^8$-sample run to $k=20$ (standard error $10^{-4}$) gives
$|a_k|=0.00495,0.00349,0.00271,0.00195,0.00142,0.00115,0.00107,0.00095,
0.00091,0.00079,0.00055,0.00045,0.00035$ for $k=8,\dots,20$, all with
sign $(-1)^k$; $k^3|a_k|$ stays in $2.4$--$2.7$ up to $k=13$ and
fluctuates in $2.8$--$3.9$ (two to three standard errors) for
$k=14$--$20$; the least-squares exponent over $8\le k\le20$ is
$2.7\pm0.3$.  From
the three-curve apart start $(0.3,0.3;\{0.4\})$: $a_1=-0.004$,
then $-0.0606,+0.0385,-0.0219,+0.0132,-0.0086,+0.0060,-0.0041,+0.0031,
-0.0020,+0.0018,-0.0014$ ($k=2,\dots,12$), with $k^3|a_k|\to2.0$--$2.4$:
the same tail, sign $(-1)^{k+B_0+1}$.  Inserting a dust curve of mass
$10^{-4}$ into either start (mass taken from an environment curve or
from an arc) reproduces every $a_k$ to within two standard errors, as
Lemma~\ref{lem:inert} predicts.  From a $9$-curve start
$(0.3,0.3;\ 7\times0.057)$, $10^8$ samples: $a_k$ passes through zero
at $k=3$ and then \emph{grows} in absolute value to $0.00145$ at
$k=7$--$8$ before decaying, with the sign $(-1)^{k+1}$ at every
$k=4,\dots,20$ (seventeen consecutive alternations); $k^3|a_k|$ rises
from $0.5$ at $k=7$ to $2$--$3$ at $k=15$--$20$.  So the tail is
universal --- $C k^{-3}$ with $C\approx2.5$--$3$ --- and is reached
after a start-dependent transient, of about eight flips when the
environment starts with nine curves.  From a $17$-curve start
($5\times10^6$ samples) $|a_k|\le7\times10^{-4}$ for $3\le k\le16$,
the transient not yet over.  From the
burnt-in ``stationary'' start of session~11 (gap window
$[\tfrac18,\tfrac38)$, $5\times10^6$ samples), $a_k(\mu)$ is at the noise
floor $4.5\times10^{-4}$ for $k\ge5$ --- but that average carries the
factor $(-1)^{B(\bar x)}$ of \eqref{eq:akB} inside the mixture, which
cancels the tail; the sign-blind $\E_\mu[a_k(\bar x)^2]$ is at
$(4\pm5)\times10^{-4}$ for $k=5..8$, consistent with the fixed-start
values.  The decomposition of \eqref{eq:akB} by $b$ from the one-curve
start (\texttt{s12\_akB}) shows the mechanism: $\Pr[B_{\tau_k}=1]$
decays like $c/k$ with $c\approx0.5$ (the chain merges back to one curve with
probability $\asymp1/t$, and there $\tau_k$ is even), the law of
$B_{\tau_k}$ spreads only like $\log k$, and the alternating sum of a
lattice law that spreads logarithmically decays polynomially.

\emph{Status: MEASURED.  Conclusions.  (i) (C2) in the geometric form
of \S\ref{sec:s11left} is false on the full state space, and on
$\Omega_K$ it holds only with $1-\rho(K)\le e^{-cK}$ (the rate the
killing provides), which is outside the budget; the session-11 value
$\rho\approx0.3$ described the transient $k\le4$.  (ii) The polynomial
decay is a $+1$ phenomenon, not a $-1$ phenomenon: the generating
function $\sum_ka_kw^k$ is singular at $w=-1$, i.e.\ at the eigenvalue
$-1$ of $\kappa_{-1}$, which is the eigenvalue $+1$ of $\kappa_1$ (the
trivial one), approached sub-geometrically by the slow variable $B$.
Nothing in the measurement bears on (C1).  (iii) The decay is
sign-alternating in $k$ from every start tested, with the sign
$(-1)^{k+B_0+1}$; this is the statement that $B_{\tau_k}\equiv k+1$ is
slightly favoured, i.e.\ that the block count at the $k$-th flip
remembers that it started odd --- a memory that $\log$-growth of $B$
erases only polynomially.}

\subsubsection{What is left, restated}\label{sec:s12left}

\begin{proposition}[polynomial parity decay suffices at $z=-1$]
\label{prop:holder}
Fix $K$ and $S'_K$.  Suppose
\begin{itemize}[leftmargin=2em]
\item[(C2$''$)] there are $C,p$ and $\delta_0>0$ with
$\sup_{\bar x\in S'_K}\bigl|\E_{\bar x}[w^{\tau_k}(-1)^{E_{\tau_k}}]\bigr|\le Ck^{-p}$
for all $k\ge1$ and all $|w|\le1$ with $|1-w|\le\delta_0$;
\item[(M$_1$)] $\sup_{\bar x\in S'_K}\E_{\bar x}[\tau_k]\le C_1k$ for all $k$.
\end{itemize}
Then $G(z)=(I+\kappa_z)^{-1}(I-\kappa_z)\mathbf 1$, evaluated on $S'_K$,
is H\"older continuous of exponent $(p-1)/(p+1)$ on
$\{|z|\le1,\ |1+z|\le\delta_0\}$, with constant
$\le C_2\,(C+C_1)$.  In particular $p>13/5$ gives an exponent $>4/9$.
\end{proposition}

\begin{proof}
By Proposition~\ref{prop:repair}, for $z=-w$,
$G(-w)=\bar M(I-\kappa_w)^{-1}(I+\kappa_w)\bar M\mathbf 1
=\mathbf 1+2\bar M\sum_{k\ge1}g_k(w)$ with
$g_k(w)(\bar x)=\E_{\bar x}[w^{\tau_k}(-1)^{E_{\tau_k}}]$
(Neumann series, $\|\kappa_w\|\le|w|<1$; the boundary values are limits).
For $|w|,|w'|\le1$, $|w^n-w'^n|\le n|w-w'|$, so
$|g_k(w)-g_k(w')|\le\min\bigl(2Ck^{-p},\ C_1k|w-w'|\bigr)$ by (C2$''$)
and (M$_1$).  Summing with the cut at $k_*=|w-w'|^{-1/(p+1)}$:
$\sum_{k\le k_*}C_1k|w-w'|+\sum_{k>k_*}2Ck^{-p}
\le C_2(C+C_1)\,|w-w'|^{(p-1)/(p+1)}$.
\end{proof}

\emph{Status: PROVED as a reduction; the hypotheses are OPEN.  Two
comments.  (a) The exponent $(p-1)/(p+1)$ is crude: it bounds
$|g_k(w)-g_k(w')|$ by the trivial Lipschitz bound and never uses that
$g_k(1)=(-1)^{E}(-1)^ka_k$ has one sign for all $k$.  If the parity
memory factorises from the time weighting,
\[
\text{(C2$'''$)}\qquad
\E_{\bar x}[w^{\tau_k}(-1)^{\tau_k}]=\E_{\bar x}[w^{\tau_k}]\,a_k(\bar x)\,(1+O(|1-w|k)),
\]
then $\sum_kg_k(w)\approx(-1)^E\sum_kCk^{-p}\E[w^{\tau_k}]$, and with
$\E[w^{\tau_k}]\approx w^{6k}$ this is $C\,\mathrm{Li}_p(w^6)$ up to
smoother terms, whose modulus of continuity at $w=1$ is
$|1-w|^{\min(1,p-1)}$ (up to a logarithm at integer $p$): the budget
then needs only $p>13/9$.  (C2$'''$) was measured (\texttt{s12\_twist}, one-curve start,
$3\times10^7$ samples, $w=0.99,0.95,0.90$ real): the exact
factorisation fails --- the twisted parity $\E[w^{\tau_k}(-1)^{\tau_k}]$
exceeds $\E[w^{\tau_k}]\,a_k$ by a factor $1.2$ at $|1-w|k=0.08$ and
$2$--$4$ at $|1-w|k\approx0.4$--$0.8$ (short paths carry more parity
memory) --- but the uniform bound (C2$''$) holds with the same
constant, $|\E[w^{\tau_k}(-1)^{\tau_k}]|\le|a_k|$ at every $k\le16$ and
every $w$ tested.  More usefully, the sum itself was measured: with
$\Sigma(w):=\sum_k(-1)^kE_x[w^{\tau_k}(-1)^{\tau_k}]$ (all terms of one
sign), $\Sigma(1)-\Sigma(w)=0.026,\ 0.085,\ 0.143$ at
$|1-w|=0.01,0.05,0.10$ (tail beyond $k=16$ estimated from $2.5k^{-3}$),
i.e.\ a modulus of continuity $\asymp|1-w|^{0.75}$ over that range ---
above the crude $(p-1)/(p+1)\approx0.5$, below the $\min(1,p-1)$ of the
factorised heuristic, and far above the $4/9$ the budget needs.  The
same on the unit circle near $w=1$ is the first numerical task of
session~13.  With the measured $p=2.7\pm0.3$ the crude route is marginal
($p>13/5$ required); the measured modulus is not.  (b) (C2$''$) at $w=1$ is exactly the polynomial
(C2$'$): $\sup_{S'_K}|a_k|\le Ck^{-p}$.  Away from $w=1$ on the circle
Theorem~\ref{thm:idleloop} gives a geometric rate; (C2$''$) asks that
the polynomial rate survive a small complex twist near $w=1$, which is
a damping of long paths and should not spoil a parity cancellation that
is produced by the block count.  This is measurable and is the first
numerical task for session~13.}

\medskip
\noindent The open conditions are now:

\medskip
\noindent\textbf{(C1)}\quad$\|(I+\kappa_1)^{-1}\|_{L^\infty(\Omega_K)}\le C(K)$ ---
equivalently $\|(I-\kappa_{-1})^{-1}\|\le C(K)$, the summability of the
signed renewal series $\sum_k\kappa_{-1}^{\,k}G$ for every bounded $G$;
(C2$'$) is the case $G=\mathbf 1$.  Unchanged, not contradicted, and no
longer reachable by any Doeblin/ergodicity argument
(Proposition~\ref{prop:nomino}, Corollary~\ref{cor:noerg}).\\
\noindent\textbf{(C2$''$)}\quad as in Proposition~\ref{prop:holder},
with $p>13/5$ and $C=C(K)$ polynomial; measured $p\approx3$ at $w=1$.

\medskip
\emph{What a proof must look like.  (i) It must be dust-blind: by
Lemma~\ref{lem:inert}, $a_k(y_\varepsilon)=a_k(y)+O(\varepsilon\,\E_y\tau_k)$
while $(-1)^{E}$ changes sign, so any argument that treats $(-1)^E$ as
a function to be contracted in a sup-type norm on configurations is
doomed, and any argument that works must use a topology in which
$y_\varepsilon\to y$.  (ii) The natural candidate --- the transport
distance between mass configurations, under the coupling of
Lemma~\ref{lem:inert} --- is not contracted: when the dust is finally
hit, the two coupled configurations differ by a macroscopic merge
(``$r+r_0$'' against ``$r+\varepsilon$, $r_0-\varepsilon$''), so a
discrepancy of size $\varepsilon$ is amplified to size $O(1)$ rather
than resolved.  A contracting metric would have to compare coarse
\emph{laws} rather than configurations.  (iii) The environment's own
randomness is needed: an adversary allowed to choose the masses gained
by the marked pair at the merge steps (whose timing it cannot control)
can bias the parity of each inter-flip time toward a target by a fixed
amount, since a gain changes the flip hazard $2p_1p_2$ by a bounded
amount at a time of its choosing among the merge steps; so no
``uniformly over environments'' statement can be true, and the proof
must use that the true gains are size-biased picks from a partition
whose count parity is invisible to them.  (iv) Given
\eqref{eq:akB}, the cleanest target is a statement about the block
count alone: the parity of $B$ at the $k$-th flip equidistributes
polynomially.  Since the total block count from one curve has the
exact law of Theorem~\ref{thm:blockcount} at fixed times, and the flip
times are a renewal-like sampling of the time axis, this may be
provable from the exact formulas by an averaging argument, at least
for the annealed environment --- which is the form the downstream
reduction (Problem~\ref{prob:survival}, whose start is
$\mu_j$) actually consumes.}
